\subsection{RELATIONS WITH REGULAR ELEMENTS OF THE LIE ALGEBRA}

PROPOSITION 6. Let V be an open subgroup of g and let exp : V → G be an exponential map defined on V.

\begin{enumerate}
\def\labelenumi{(\roman{enumi})}
\item
  There exists a neighbourhood \(\mathrm{W}\) of 0 in \(\mathrm{V}\) such that \(\mathfrak{g}^1(\exp x) = \mathfrak{g}^0(x)\) for all \(x \in \mathrm{W}\) .
\item
  If \(k = \mathbf{R}\) or \(\mathbf{C}\) , \(\mathfrak{g}^1(\exp x) \supset \mathfrak{g}^0(x)\) for all \(x \in \mathfrak{g}\) .
\end{enumerate}

By Cor. 3 of Prop. 8 of Chap. III, §4, no. 4, there exists a neighbourhood \(V'\) of 0 in \(V\) such that, for all \(x \in V'\) , \(\exp(\mathrm{ad}(x)) = \sum_{n=0}^{\infty} \frac{1}{n!} \mathrm{ad}(x)^n\) is defined and \(\mathrm{Ad}(\exp(x)) = \exp(\mathrm{ad}(x))\) . If \(P \in k[X]\) and \(\alpha \in \mathrm{End}(\mathfrak{g})\) , it is easy to check that \(\mathfrak{g}^\lambda(\alpha) \subset \mathfrak{g}^{P(\lambda)}(P(\alpha))\) for all \(\lambda \in k\) . Consequently,

\[
\mathfrak {g} ^ {0} (\operatorname{ad} (x)) \subset \mathfrak {g} ^ {1} (\exp (\operatorname{ad} (x))) = \mathfrak {g} ^ {1} (\operatorname{Ad} (\exp x)) = \mathfrak {g} ^ {1} (\exp x)
\]

for all \(x \in V'\) . If k = R or C, V = g and we can take \(V' = V\) , which proves (ii). We prove (i). Let U be a neighbourhood of 0 in \(\text{End}(\mathfrak{g})\) such that \(\text{Log}(1 + \alpha) = \sum_{n > 0} (-1)^{n+1} \frac{1}{n} \alpha^n\) is defined for all \(\alpha \in U\) . Then \(\log \circ \exp = 1\) on a neighbourhood of 0 and \(\mathfrak{g}^1(1 + \alpha) \subset \mathfrak{g}^0 (\text{Log}(1 + \alpha))\) for all \(\alpha \in U\) . Let W be the neighbourhood of 0 in g consisting of those \(x \in V'\) such that \(\exp x \in 1 + U\) and

\[
\operatorname{Log} (\exp (\operatorname{ad} (x))) = \operatorname{ad} (x).
\]

Then, for all \(x \in W\) ,

\[
\begin{array}{c} \mathfrak {g} ^ {1} (\exp x) = \mathfrak {g} ^ {1} (\mathrm{Ad} (\exp x)) = \mathfrak {g} ^ {1} (\exp (\mathrm{ad} (x))) \\ \qquad \subset \mathfrak {g} ^ {0} (\operatorname{Log} (\exp (\operatorname{ad} (x)))) = \mathfrak {g} ^ {0} (\operatorname{ad} (x)) = \mathfrak {g} ^ {0} (x). \end{array}
\]

This shows that \(\mathfrak{g}^1 (\exp x) = \mathfrak{g}^0 (x)\) for all \(x\in \mathrm{W}\) .

Lemma 4. Let U be a neighbourhood of 0 in g and exp an exponential map from U to G, étale at every point of U and such that \(\mathfrak{g}^{1}(\exp x) = \mathfrak{g}^{0}(x)\) for all \(x \in U\) .

\begin{enumerate}
\def\labelenumi{(\roman{enumi})}
\item
  The function \(r_{\mathrm{Ad}}^0\) is constant and equal to the rank of \(\mathfrak{g}\) on \(\exp(\mathrm{U})\) .
\item
  If \(x \in \mathrm{U}\) , exp \(x\) is regular if and only if \(x\) is a regular element of \(\mathfrak{g}\) .
\item
  An element \(a \in \exp(\mathrm{U})\) is regular if and only if \(\mathfrak{g}^1(a)\) is a Cartan subalgebra of \(\mathfrak{g}\) .
\end{enumerate}

Let \(l = \operatorname{rk}(\mathfrak{g})\) . If \(x \in \mathrm{U}\) is a regular element of \(\mathfrak{g}\) ,

\[
r _ {\mathrm{Ad}} (\exp x) = \dim \mathfrak {g} ^ {1} (\exp x) = \dim \mathfrak {g} ^ {0} (x) = l.
\]

Since the regular elements of g belonging to U constitute a neighbourhood of x and exp is étale at x, this shows that exp x is regular and that \(r_{\mathrm{Ad}}^{0}(\exp x) = l\) . The regular elements of g belonging to U being dense in U, we have \(r_{\mathrm{Ad}}^{0}(a) = l\) for all \(a \in \exp(\mathrm{U})\) . Let \(a \in \exp(\mathrm{U})\) be a regular element of G and let \(x \in U\) be such that \(a = \exp x\) . Since \(\mathfrak{g}^{0}(x) = \mathfrak{g}^{1}(a)\) , \(\dim \mathfrak{g}^{0}(x) = r_{\mathrm{Ad}}^{0}(a) = l\) . Consequently, x is a regular element of g and \(\mathfrak{g}^{1}(a)\) is a Cartan subalgebra of g. Finally, if \(a \in \exp(\mathrm{U})\) and \(\mathfrak{g}^{1}(a)\) is a Cartan subalgebra of g,

\[
r _ {\mathrm{Ad}} (a) = \dim \mathfrak {g} ^ {1} (a) = l = r _ {\mathrm{Ad}} ^ {0} (a),
\]

so \(a\) is regular.

PROPOSITION 7. Let V be a neighbourhood of e in G. Every Cartan sub-algebra of g is of the form \(\mathfrak{g}^{1}(a)\) where a is a regular element of G belonging to V.

By Prop. 6, there exists an open neighbourhood U of 0 in g and an exponential map \(\exp: U \to G\) satisfying the conditions of Lemma 4. If h is a Cartan subalgebra of g, there exists a regular element \(x \in h\) such that \(\mathfrak{h} = \mathfrak{g}^{0}(x)\) (§3, Th. 2). On the other hand, there exists an element \(t \in k^{*}\) such that \(tx \in U\) and \(\exp(tx) \in V\) . Then \(\mathfrak{h} = \mathfrak{g}^{0}(x) = \mathfrak{g}^{0}(tx) = \mathfrak{g}^{1}(\exp(tx))\) , and by Lemma 4 (ii), \(\exp(tx)\) is a regular element of G.

PROPOSITION 8. Let \(l\) be the rank of \(\mathfrak{g}\) . There exists an open subgroup \(G^*\) of \(G\) such that:

\begin{enumerate}
\def\labelenumi{(\roman{enumi})}
\item
  the function \(r_{\mathrm{Ad}}^0\) is constant on \(\mathbf{G}^*\) and its value is \(l\) ;
\item
  an element \(a \in \mathbf{G}^*\) is regular if and only if \(\mathfrak{g}^1(a)\) is a Cartan subalgebra of \(\mathfrak{g}\) ;
\item
  if \(a \in \mathbf{G}^*\) , every Cartan subalgebra of \(\mathfrak{g}^1(a)\) is a Cartan subalgebra of \(\mathfrak{g}\) .
\item
  By Prop. 6, there exists an open neighbourhood U of 0 in g and an exponential map exp from U to G satisfying the conditions of Lemma 4. In what follows, \(G^{*}\) will denote the identity component of G if k = R or C and an open subgroup of G contained in \(\exp(\mathrm{U})\) if k is ultrametric. Since \(r_{Ad}^{0}\) is locally constant and its value at any point of \(\exp(\mathrm{U})\) is l (Lemma 4 (i)), it follows that \(r_{Ad}^{0}\) is constant and equal to l on \(G^{*}\) .
\item
  Let \(\mathrm{R}^*\) (resp. \(\mathrm{S}^*\) ) be the set of regular elements of \(\mathrm{G}^*\) (resp. the set of elements \(a \in \mathrm{G}^*\) such that \(\mathfrak{g}^1(a)\) is a Cartan subalgebra of \(\mathfrak{g}\) ). Then \(\mathrm{S}^* \subset \mathrm{R}^*\) . Indeed, if \(a \in \mathrm{S}^*\) , then \(r_{\mathrm{Ad}}(a) = l = r_{\mathrm{Ad}}^0(a)\) . We show that \(\mathrm{R}^* \subset \mathrm{S}^*\) . If \(k\) is ultrametric, this follows from the inclusion \(G^{*} \subset \exp(U)\) and Lemma 4 (iii). Assume that k = C. By the Cor. of Prop. 5, if \(a \in R^{*}\) , then for every \(a'\) belonging to a neighbourhood of a, \(\mathfrak{g}^{1}(a')\) is conjugate to \(\mathfrak{g}^{1}(a)\) by an automorphism of g. This proves that \(S^{*}\) and \(R^{*} - S^{*}\) are open subsets of \(G^{*}\) . We have seen that \(S^{*}\) contains all the regular elements in a neighbourhood of e (Lemma 4 (iii)); consequently, \(S^{*}\) is non-empty. Since \(G^{*}\) is connected, so is \(R^{*}\) (Prop. 1) and consequently \(S^{*} = R^{*}\) .
\end{enumerate}

It remains to study the case \(k = \mathbf{R}\) . Assume first of all that \(\mathbf{G}^*\) is an integral subgroup of \(\mathbf{GL}(\mathrm{E})\) where \(\mathrm{E}\) denotes a finite dimensional real vector space. Let \(\mathrm{G}_c^*\) be the integral subgroup of \(\mathbf{GL}(\mathrm{E} \otimes_{\mathbf{R}} \mathbf{C})\) with Lie algebra \(\mathfrak{g}_c = \mathfrak{g} \otimes \mathbf{C}\) . There exists an analytic function on \(\mathrm{G}_c^*\) whose set of zeros is the complement of the open set of regular elements of \(\mathrm{G}_c^*\) . By Differentiable and Analytic Manifolds, Results, 3.2.5, this function cannot vanish at every point of \(\mathrm{G}^*\) . Consequently, \(\mathrm{G}^*\) contains a regular element of \(\mathrm{G}_c^*\) . Let \(\mathrm{Ad}_c\) be the adjoint representation of \(\mathrm{G}_c^*\) . For any \(a \in \mathrm{G}^*\) , \(\mathfrak{g}_c^1(a) = \mathfrak{g}^1(a) \otimes \mathbf{C}\) , so \(r_{\mathrm{Ad}_c}(a) = r_{\mathrm{Ad}}(a)\) . If \(a \in \mathrm{G}^*\) is a regular element of \(\mathrm{G}_c^*\) , this is a regular element of \(\mathrm{G}^*\) and \(r_{\mathrm{Ad}_c}^0(a) = r_{\mathrm{Ad}}^0(a)\) . The functions \(r_{\mathrm{Ad}_c}^0\) and \(r_{\mathrm{Ad}}^0\) being constant on \(\mathrm{G}_c^*\) and on \(\mathrm{G}^*\) , respectively, it follows that the regular elements of \(\mathrm{G}^*\) are the regular elements of \(\mathrm{G}_c^*\) belonging to \(\mathrm{G}^*\) . From the above, if \(a\) is a regular element of \(\mathrm{G}^*\) , \(\mathfrak{g}_c^1(a) = \mathfrak{g}^1(a) \otimes \mathbf{C}\) is a Cartan subalgebra of \(\mathfrak{g}_c\) ; this implies that \(\mathfrak{g}^1(a)\) is a Cartan subalgebra of \(\mathfrak{g}\) (§2, Prop. 3).

Assume now that G is simply connected. There exists a finite dimensional real vector space E and an étale morphism f from G to an integral subgroup \(G'\) of \(\mathbf{GL}(\mathrm{E})\) (Chap. III, §6, no. 1, Cor. of Th. 1). By Prop. 2, if \(a \in G\) is regular, \(f(a)\) is regular. By the preceding, \(\mathfrak{g}'^{1}(f(a))\) is a Cartan subalgebra of the Lie algebra \(g'\) of \(G'\) . Since \(\mathfrak{g}'^{1}(f(a)) = (\mathrm{T}f)\mathfrak{g}^{1}(a)\) and Tf is an isomorphism from g to \(g'\) , this proves that \(\mathfrak{g}^{1}(a)\) is a Cartan subalgebra of g.

We turn finally to the general case \((k = \mathbf{R})\) . Let \(\tilde{G}\) be a universal covering of \(G^{*}\) , \(\tilde{\mathfrak{g}} = \mathrm{L}(\tilde{\mathrm{G}})\) , and q the canonical map from \(\tilde{G}\) to \(G^{*}\) . Since the kernel of q is contained in the centre of \(\tilde{G}\) , if \(a \in G^{*}\) is regular and if \(a' \in q^{-1}(a)\) , then \(a'\) is regular (Prop. 2). By the preceding, \(\tilde{\mathfrak{g}}^{1}(a')\) is a Cartan subalgebra of \(\tilde{g}\) . Since \(\mathfrak{g}^{1}(a) = (\mathrm{T}q)\tilde{\mathfrak{g}}^{1}(a')\) and since Tq is an isomorphism from \(\tilde{g}\) to g, this proves that \(\mathfrak{g}^{1}(a)\) is a Cartan subalgebra of g.

\begin{enumerate}
\def\labelenumi{(\roman{enumi})}
\setcounter{enumi}{2}
\tightlist
\item
  By Prop. 5, there exists a neighbourhood \(\mathrm{V}\) of \(a\) such that, for all \(a' \in \mathrm{V}\) , \(\mathfrak{g}^1(a')\) is conjugate to a subalgebra of \(\mathfrak{g}^1(a)\) by an automorphism of \(\mathfrak{g}\) . Since every neighbourhood of \(a\) contains a regular element of \(\mathbf{G}^*\) , it follows from (ii) that \(\mathfrak{g}^1(a)\) contains a Cartan subalgebra of \(\mathfrak{g}\) . Thus, by Prop. 3 of §3, every Cartan subalgebra of \(\mathfrak{g}^1(a)\) is a Cartan subalgebra of \(\mathfrak{g}\) .
\end{enumerate}

Remark. If \(k = \mathbf{C}\) , the subalgebras \(\mathfrak{g}^1(a)\) , for \(a\) regular and belonging to a connected component \(M\) of \(G\) , are conjugate under \(\operatorname{Int}(\mathfrak{g})\) . Indeed, let \(R\) be the set of regular elements of \(G\) . For all \(a \in R \cap M\) , let \(M_a\) be the set of those \(b \in R \cap M\) such that \(\mathfrak{g}^1(a)\) is conjugate to \(\mathfrak{g}^1(a)\) under \(\operatorname{Int}(\mathfrak{g})\) . We have \(\operatorname{Int}(\mathfrak{g}) = \operatorname{Ad}(G^0)\) , where \(G^0\) is the identity component of \(G\) . By the Corollary to Prop. 5, \(M_{a}\) is open in R. It follows that \(M_{a}\) is open and closed in R. Since \(k = C\) , \(R \cap M\) is connected (Lemma 1), hence \(M_{a} = R \cap M\) .

